\phantomsection
\subsection*{M4C. Polish Mead}
\addcontentsline{toc}{subsection}{M4C. Polish Mead}

\textit{Polish Mead é uma família de estilos com origem e caráter em comum, definida legalmente na União Europeia (UE) como Especialidade Tradicional Garantida, sob o nome \textbf{Staropolski Tradycyjny} (Polonês Histórico Tradicional). A definição de estilo do BJCP está alinhada a essas especificações da UE. O Polish Mead é classificado de acordo com a proporção entre os volumes de mel e água utilizados, com as seguintes variantes: \textbf{Czwórniak} (1:3, aproximadamente equivalente a um hidromel Médio), \textbf{Trójniak} (1:2, aproximadamente equivalente a um hidromel Doce), \textbf{Dwójniak} (1:1) e \textbf{Półtorak} (1:0,5).}

\textit{No entanto, sem um caráter distinto, o Polish Mead seria apenas mais um hidromel produzido segundo proporções rígidas de receita. Portanto, este estilo descreve uma parte da família dos hidroméis poloneses que apresenta um caráter tradicional distinto de maturação e sensação na boca.}

\textbf{Impressão Geral}: As versões mais fortes ou densas de Polish Mead são caracterizadas por uma sensação na boca viscosa, caráter maturado e, frequentemente, final doce, enquanto as versões menos alcoólicas ou menos densas podem parecer, superficialmente, semelhantes aos hidroméis da categoria M1 Traditional Mead e às suas variações com frutas ou especiarias. Este estilo destina-se a hidroméis que apresentem esse caráter polonês distinto.

\textbf{Aroma}: O caráter e a intensidade do mel variam de acordo com as variedades utilizadas e com o teor alcoólico e o dulçor do hidromel. Um caráter maturado é comum e pode apresentar notas positivas de oxidação, de intensidade leve a moderada, resultantes naturalmente da maturação. O nível de dulçor depende da densidade final, mas pode ser elevado em hidroméis produzidos com maior proporção de mel.

Qualquer adição de frutas deve apresentar o mesmo caráter maturado do mel, e não um caráter fresco ou suculento, podendo variar de intensidade muito baixa a média. Adições de especiarias ou ervas variam de médio-baixas a ausentes e não devem dominar o mel ou as frutas. Alguma presença de álcool pode ser perceptível no aroma, mas a maturação e o dulçor frequentemente mascaram ou suavizam esse caráter.

\textbf{Aparência}: Límpido. Still. Quando jovem, a cor varia de dourado a âmbar escuro, dependendo das variedades de mel utilizadas. Com a maturação e a adição de frutas, a cor pode se tornar mais escura, menos vívida ou desenvolver nuances amarronzadas.

\textbf{Sabor}: Semelhante ao aroma quanto ao caráter e à intensidade do mel, às adições de frutas e especiarias, ao álcool e ao caráter de maturação. O final se torna mais doce nas versões com maior proporção de mel em relação à água, mas não deve apresentar sabor de mel não fermentado nem ser enjoativo.

A acidez e os taninos, às vezes provenientes das adições de frutas e especiarias, podem moderar essa percepção de dulçor, assim como uma maturação prolongada. Os níveis de acidez e taninos frequentemente têm menor destaque no equilíbrio do que em muitos outros estilos de hidromel. Nos hidroméis com frutas ou especiarias, o equilíbrio normalmente favorece o mel, com os ingredientes adicionados desempenhando um papel complementar, sem dominar, especialmente nos exemplares bem maturados.

A maturação frequentemente se manifesta por meio da maciez, suavidade e integração dos sabores, incluindo a moderação de um dulçor intenso. Também pode ser percebida pelas notas positivas encontradas em exemplares envelhecidos de Sherry, Porto, Marsala e Madeira, como notas de nozes ou caráter de frutas escuras ou secas. Entre os defeitos estão sabores acéticos ou de solvente, acidez pungente, aspereza e oxidação excessiva.

\textbf{Sensação na Boca}: Carbonatação Still. As versões mais densas apresentam corpo viscoso, mas não devem ser enjoativas nem excessivamente pesadas. A maturação pode proporcionar uma textura aveludada, sedosa e macia. O corpo varia de médio-leve a cheio em um \textit{Czwórniak}; de médio a cheio em um \textit{Trójniak}; de médio-alto a viscoso em um \textit{Dwójniak}; e de cheio a viscoso em um \textit{Półtorak}. A percepção de aquecimento alcoólico varia conforme a versão. Versões com frutas, especiarias, lúpulo ou madeira podem apresentar níveis mais elevados de taninos ou acidez, que compensam parcialmente o dulçor e a viscosidade.

\textbf{Ingredientes}: Tradicionalmente produzido com méis da região, como tília, trigo-sarraceno, acácia ou diversos méis de campo ou floresta. Pode conter sucos de frutas, especiarias ou lúpulo. Entre as frutas locais estão maçã, pera, ameixa, damasco, cereja, framboesa, groselha-preta ou vermelha, amora, mirtilo, lingonberry, groselha-espinhosa, uvas ou frutas silvestres da floresta. As especiarias podem incluir aquelas utilizadas em confeitaria ou especiarias típicas do outono e inverno da região. Um caráter discreto de madeira (carvalho) é aceitável, mas não obrigatório.

\textbf{Comentários}: Produto tradicional elaborado há mais de um milênio em uma região que abrange os atuais territórios da Polônia e da Lituânia, além de partes da Bielorrússia e da Ucrânia. Ingredientes provenientes dessa região são considerados tradicionais. Todas as variantes de Polish Mead podem conter frutas, especiarias e lúpulo. Quando frutas são utilizadas, o uso de suco é tradicional, substituindo pelo menos 30\% do volume de água nas proporções estabelecidas. \textit{Dwójniak} e \textit{Półtorak} são tradicionalmente produzidos inicialmente como \textit{Trójniak}, recebendo o restante do mel próximo ao final da fermentação ou durante a maturação. Versões modernas podem ser fortificadas com destilados. Versões modernas que utilizem ingredientes não tradicionais ou outras proporções de mel e água devem ser inscritas como M4F Experimental Mead, desde que mantenham o caráter polonês.

\textbf{Estatísticas}: Os valores de OG referem-se a hidroméis sem adição de frutas; exemplares com frutas apresentam valores mais elevados.

\begin{center}
\footnotesize
\begin{tabular}{lcccc}
\toprule
Variante & ABV & Maturação Mín. & OG & FG \\
\midrule
Czwórniak & 9-12\% & 9 meses & 1.090 & 1.012-1.033 \\
Trójniak & 12-15\% & 1 ano & 1.144 & 1.024-1.045 \\
Dwójniak & 15-18\% & 2 anos & 1.215 & 1.065-1.089 \\
Półtorak & 15-18\% & 3 anos & 1.286 & 1.150-1.170 \\
\bottomrule
\end{tabular}
\end{center}

\textbf{Instruções para Inscrição}: Os participantes devem especificar a variante de Polish Mead (\textit{Czwórniak}, \textit{Trójniak}, \textit{Dwójniak} ou \textit{Półtorak}). Presume-se que o teor alcoólico e a carbonatação estejam dentro das especificações. O nível de dulçor deve ser declarado. Quaisquer ingredientes especiais (frutas, especiarias, lúpulo) devem ser declarados. Variedades de mel podem ser declaradas.

\textbf{Exemplos Comerciais}: Ami Honey Trójniak Sandomierski, Apis Czwórniak Korzenny, Apis Półtorak Jadwiga, Augustowska Trójniak Jabłko na Miodzie, Dziki Miód Trójniak Porzeczkowiec, Jaros Dwójniak Maliniak, Jaros Trójniak Trybunalski, Mazurskie Półtorak Piasecki